\input{../tex-inputs/latex-leseansicht-vorspann}

\section[Gustav und Max Schwarzkopf an Arthur Schnitzler, 14. 7. 1906]{L04273 Gustav und Max Schwarzkopf an Arthur Schnitzler, 14. 7. 1906}
\nopagebreak\mylabel{L04273v}
\rehead{ }\normalsize\beginnumbering\briefempfaengerindex{Schnitzler, Arthur@\textsc{Schnitzler, Arthur}!zzzSchwarzkopf, Max@\emph{von Max Schwarzkopf}!1906-07-142@{14. 7. 1906}|(be}\briefempfaengerindex{Schnitzler, Arthur@\textsc{Schnitzler, Arthur}!zzzSchwarzkopf, Gustav@\emph{von Gustav Schwarzkopf}!1906-07-142@{14. 7. 1906}|(be}
\toendnotes[C]{\smallbreak\pagebreak[2]}
\correspDesc{Versand  durch Gustav Schwarzkopf, Max Schwarzkopf am 14. 7. 1906 in Wien
\newline{}Übermittlung  am 14. 7. 1906 in Wien
\newline{}Weiterleitung  am 16. 7. 06 in Helsingør
\newline{}Erhalt  durch Arthur Schnitzler im Zeitraum [15. 7. 1906
                  – 19. 7. 1906?] in Marienlyst}\toendnotes[C]{\smallbreak}
\Standort{CUL, Schnitzler, B 96a.}
\physDesc{Postkarte, 843 Zeichen
\newline{}Handschrift Gustav Schwarzkopf: schwarze Tinte, deutsche Kurrent
\newline{}Handschrift Max Schwarzkopf: schwarze Tinte, deutsche Kurrent
\newline{}Versand: 1) Stempel: »\nobreak{}\oindex{I., Innere Stadt@\textbf{I., Innere Stadt}, \emph{Verwaltungsgebiet}|pwk}1/1 Wien 1, 14. 7. 06, N 7–8\nobreak{}«.   2) Stempel: »\nobreak{}\oindex{Helsingør@\textbf{Helsingør}, \emph{Hauptstadt}|pwk}Helsingør, 16 7 06, 12F\nobreak{}«. }\toendnotes[C]{\smallbreak}\pstart{}{\pb}I. Tiefer Graben 23. II. St. Th. 6\oindex{Wien@\textbf{Wien}!I., Innere Stadt@\textbf{I., Innere Stadt}!Tiefer Graben 23@\textbf{Tiefer Graben 23}, \emph{Wohngebäude}|pw}\pend{}\pstart{}Wien\oindex{Wien@\textbf{Wien}, \emph{Verwaltungsgebiet}|pw}\pend{}{\bigskip}\pstart{}Herrn \textsc{D\textsuperscript{r} Arthur
                     Schnitzler}\pend{}\pstart{}\textsc{Marienlyst\oindex{Marienlyst@\textbf{Marienlyst}, \emph{Gut}|pw}}\pend{}\pstart{}\textsc{(Dänemark\oindex{Dänemark@\textbf{Dänemark}|pw}) Kurhôtel\oindex{Kurhotellet@\textbf{Kurhotellet}, \emph{Hotel}|pw}}\pend{}{\bigskip}\vspace{1em}
\pstart
           \raggedleft{}{\pb}14. VII 06\pend
           \vspace{0.5em}
\pstart
           Lieber Arthur, wie Sie{ }ſehen bin ich noch in Wien\oindex{Wien@\textbf{Wien}, \emph{Verwaltungsgebiet}|pw}. An »lockenden Anerbietungen« fehlt es nicht, aber ich ko{\geminationn}te mich noch nicht entſchließen – das Wetter iſt
               allerdings{ }ſehr rückſichtsvoll für die 1 ½ Millionen, die noch hier{ }ſind, quält
               uns nicht mit »ſengender Hitze«; heute iſt es{ }ſogar{ }ſo kalt, daß man eher geneigt iſt
               heizen zu laſſen, als eine So{\geminationm}erfriſche aufzuſuchen. Von
                  Beka{\geminationn}ten iſt freilich Niemand mehr hier,{ }ſelbſt Eberma{\geminationn}\pwindex{Ebermann, Leo 16.\,7.\,1863 Draganovka – 9.\,10.\,1914 Wien@\textsc{Ebermann, Leo} (16.\,7.\,1863 Draganovka – 9.\,10.\,1914 Wien), \emph{Schriftsteller, Journalist, Rechtswissenschaftler}|pw} iſt fort. In einem Neſt\oindex{?? [Urlaubsort von Leo Ebermann, Juli 1906]@\textbf{?? [Urlaubsort von Leo Ebermann, Juli 1906]}|pwv} an der Aspang-Bahn\orgindex{Aspangbahn@Aspangbahn|pw}, zur Beruhigung{ }ſeiner Nerven. Ich habe ihn einmal beſucht. So{\geminationn}tag war
               ich in Rodaun\oindex{Wien@\textbf{Wien}!XXIII., Liesing@\textbf{XXIII., Liesing}!Rodaun@\textbf{Rodaun}, \emph{Region}|pw} bei Richard\pwindex{Beer-Hofmann, Richard 11.\,7.\,1866 Wien – 26.\,9.\,1945 New York City@\textsc{Beer-Hofmann, Richard} (11.\,7.\,1866 Wien – 26.\,9.\,1945 New York City), \emph{Schriftsteller}|pw}. Wiſſen Sie vielleicht \strikeout{ob} wie der \label{K_L04273-1v}\edtext{Prozeß des Paul M.\pwindex{Marx, Paul 21.\,7.\,1879 Wien – 30.\,10.\,1956 ebd.@\textsc{Marx, Paul} (21.\,7.\,1879 Wien – 30.\,10.\,1956 ebd.), \emph{Regisseur, Schauspieler}|pw} mit Düſſeldorf\oindex{Düsseldorf@\textbf{Düsseldorf}|pw}\orgindex{Schauspielhaus Düsseldorf@Schauspielhaus Düsseldorf|pwv}}{\lemma{\textnormal{\emph{Prozeß … Düsseldorf}}}\Cendnote{\textnormal{Paul Marx\pwindex{Marx, Paul 21.\,7.\,1879 Wien – 30.\,10.\,1956 ebd.@\textsc{Marx, Paul} (21.\,7.\,1879 Wien – 30.\,10.\,1956 ebd.), \emph{Regisseur, Schauspieler}|pwk} war eine Saison lang am \emph{Schauspielhaus Düsseldorf}\orgindex{Schauspielhaus Düsseldorf@Schauspielhaus Düsseldorf|pwk} engagiert gewesen. Für den 18. 5. 1906 notierte
                  sich Schnitzler im \emph{Tagebuch}\pwindex{Schnitzler, Arthur 15.\,5.\,1862 Wien – 21.\,10.\,1931 ebd.@\textsc{Schnitzler, Arthur} (15.\,5.\,1862 Wien – 21.\,10.\,1931 ebd.), \emph{Schriftsteller, Mediziner}!Tagebuch@\strich\emph{Tagebuch}|pwk}: »Paul
                        M.\pwindex{Marx, Paul 21.\,7.\,1879 Wien – 30.\,10.\,1956 ebd.@\textsc{Marx, Paul} (21.\,7.\,1879 Wien – 30.\,10.\,1956 ebd.), \emph{Regisseur, Schauspieler}|pw}, nach Düsseldorfer\oindex{Düsseldorf@\textbf{Düsseldorf}|pw} Aerger (Schauspielhaus\orgindex{Schauspielhaus Düsseldorf@Schauspielhaus Düsseldorf|pw}{ }Dumont\pwindex{Dumont, Louise 22.\,2.\,1862 Köln – 16.\,5.\,1932 Düsseldorf@\textsc{Dumont, Louise} (22.\,2.\,1862 Köln – 16.\,5.\,1932 Düsseldorf), \emph{Theaterleiterin, Schauspielerin}|pw} – Lindemann\pwindex{Lindemann, Gustav 24.\,8.\,1872 Danzig – 6.\,5.\,1960 Sonnenholz@\textsc{Lindemann, Gustav} (24.\,8.\,1872 Danzig – 6.\,5.\,1960 Sonnenholz), \emph{Theaterleiter, Schauspieler}|pw}) wieder in Wien\oindex{Wien@\textbf{Wien}, \emph{Verwaltungsgebiet}|pw}. (Zu Brahm\pwindex{Brahm, Otto 5.\,2.\,1856 Hamburg – 28.\,11.\,1912 Berlin@\textsc{Brahm, Otto} (5.\,2.\,1856 Hamburg – 28.\,11.\,1912 Berlin), \emph{Theaterleiter, Regisseur}|pw}
                     engagirt.)« Über einen etwaigen Prozess ist nichts Näheres bekannt.}}}\label{K_L04273-1} ausgefallen iſt? Empfehlen Sie mich Ihrer Frau\pwindex{Schnitzler, Olga 17.\,1.\,1882 Wien – 13.\,1.\,1970 Lugano@\textsc{Schnitzler, Olga} (17.\,1.\,1882 Wien – 13.\,1.\,1970 Lugano), \emph{Schauspielerin, Sängerin}|pwv} und bringen Sie mich bei
                  Heini\pwindex{Schnitzler, Heinrich 9.\,8.\,1902 Hinterbrühl – 12.\,7.\,1982 Wien@\textsc{Schnitzler, Heinrich} (9.\,8.\,1902 Hinterbrühl – 12.\,7.\,1982 Wien), \emph{Regisseur, Schauspieler}|pw} in Eri{\geminationn}erung.\pend
           
\pstart
           Ihr{\\[\baselineskip]}\spacefill\mbox{Gustav Sch.}\pend
           \leftskip=0em{}\selectlanguage{ngerman}\vspace{1em}{\vspace{1\baselineskip}}
\pstart
           {[}hs. Schwarzkopf:{]} \label{T_L04273-1v}\edtext{Beſten
                  Gruß}{\lemma{\textnormal{\emph{Besten
                  Gruß}}}\Cendnote{\textnormal{ab hier seitlich am linken
                  Rand}}}\label{T_L04273-1}\pend
           \pstart \spacefill\mbox{Max}\pend{}\selectlanguage{ngerman}\endnumbering\briefempfaengerindex{Schnitzler, Arthur@\textsc{Schnitzler, Arthur}!zzzSchwarzkopf, Max@\emph{von Max Schwarzkopf}!1906-07-142@{14. 7. 1906}|)be}\briefempfaengerindex{Schnitzler, Arthur@\textsc{Schnitzler, Arthur}!zzzSchwarzkopf, Gustav@\emph{von Gustav Schwarzkopf}!1906-07-142@{14. 7. 1906}|)be}\mylabel{L04273h}
\begin{anhang}
\end{anhang}\newcommand{\dateiname}{L04273}\newcommand{\titel}{Gustav und Max Schwarzkopf an Arthur Schnitzler, 14. 7. 1906}\newcommand{\editorInnen}{Herausgegeben von Jahnke, SelmaMüller, Martin Anton}\input{../tex-inputs/latex-leseansicht-abspann}
